\centerline{\Huge \bf  Домашнее задание }
\centerline{\Huge \bf Сумма цифр числа}

\begin{flushright}
    \scriptsize{Чтобы найти решение, найти которое невозможно, \\
                нужно допускать ошибки, иначе все бы знали, как оно решается\\
                }
\end{flushright}

\problem{1!} Обозначим через $S(x)$ сумму цифр натурального числа x. Решите уравнение:
$x + S(x) + S(S(x)) + S(S(S(x))) = 2048$

\problem{2!} АА написал на доске натуральные числа $N$ и $N+1$. Он заметил, что сумма цифр каждого из этих чисел делится на $26$. Оказалось, что $N$ - наименьшее число, для которого это условие выполняется. Найдите сумму цифр числа N и само число $N$.

\problem{3} Обозначим через S(x) сумму цифр натурального числа x. Решите уравнение:
$x + S(x) + S(S(x))  = 7777$

\problem{4} АА написал на доске натуральные числа $N$  и  $N+3$. Он заметил, что сумма цифр каждого из этих чисел делится на $9$. Оказалось, что $N$ - наименьшее число, для которого это условие выполняется. Найдите сумму цифр числа $N$ и само число $N$.